% Generated draft: requires mathematical and source review.
\documentclass[11pt]{article}
\usepackage[margin=25mm]{geometry}
\usepackage{fontspec}
\setmainfont{Latin Modern Roman}
\usepackage{amsmath,amssymb,mathtools,unicode-math}
\setmathfont{Latin Modern Math}
\usepackage{enumitem,needspace,xurl,hyperref}
\setlength{\parindent}{0pt}
\setlength{\parskip}{5pt}
\setlength{\emergencystretch}{3em}
\newcommand{\N}{\mathbb{N}}
\newcommand{\Z}{\mathbb{Z}}
\newcommand{\Q}{\mathbb{Q}}
\newcommand{\R}{\mathbb{R}}
\newcommand{\C}{\mathbb{C}}
\newcommand{\F}{\mathbb{F}}
\newcommand{\A}{\mathbb{A}}
\newcommand{\PP}{\mathbb{P}}
\newcommand{\E}{\mathbb{E}}
\newcommand{\eps}{\varepsilon}
\newcommand{\dd}{\,\mathrm{d}}
\begin{document}
\section*{UnsolvedMath ID: 1900010}
\textbf{Source problem code:} catalogue entry
\subsection{Definitions and mathematical statement}
A standard Klein-sail construction starts with a pointed cone \(C\subseteq\R^d\), with vertex \(0\), and takes the polyhedral boundary facing \(0\) of
\[
\operatorname{conv}\bigl((C\cap\Z^d)\setminus\{0\}\bigr).
\]
Write this sail as \(S(C)\), with its face decomposition. When a specified group \(U_C\) of integer-linear symmetries preserves the sail and acts freely and properly with compact quotient homeomorphic to a torus, the quotient faces give a torus decomposition \(D(C)=S(C)/U_C\). Periodicity, the dimension, and the choice of the quotient symmetry group are required to make this construction definite; none is specified in the excerpt.

Arnold's supplied conjecture is that integer-noncongruent Klein sails have distinct torus decompositions. Once the intended admissible class, integer congruence relation \(\cong_{\Z}\), and equivalence relation \(\cong_D\) on torus decompositions are fixed, its precise injectivity form is
\[
\forall S,S'\text{ in that class},\qquad
D(S)\cong_D D(S')\ \Longrightarrow\ S\cong_{\Z}S'.
\]
For example, a possible integer-linear convention is \(S'=A(S)\) for \(A\in\operatorname{GL}_d(\Z)\), but the source excerpt does not distinguish this from orientation-preserving or integer-affine congruence. Likewise, ``distinct decompositions'' could compare unlabelled cell structures or cell structures carrying lattice-geometric labels and markings; those alternatives are not interchangeable. The displayed implication is consequently the exact logical target conditional on those missing definitions, rather than a selection of one interpretation.

The supplied moderate-confidence OPEN-TRIAGE summary treats the general conjecture as open, says no proof or counterexample was found in its retrieved literature, and mentions discussions in Arnold's problems and continued-fraction work involving Markov spectra and cubic fields. That search assessment is not an independently verified literature conclusion.
\subsection{Short English statement}
Does the torus decomposition determine a periodic Klein sail up to integer congruence? The allowed sails, quotient group, congruence convention, and information retained in the decomposition must be specified.
\subsection{Sources}
\begin{itemize}
\item Local UnsolvedMath registry, permanent ID \(1900010\), exact source code \texttt{catalogue entry}, ``Geometry of Continued Fractions --- Combinatorial structure of sails.''
\item Karpenkov, \emph{Open problems in geometry of continued fractions} (2017), Conjecture 10, attributed to V. Arnold; \url{https://arxiv.org/abs/1712.01450}. Source-TeX extraction; recorded access 2026-07-29.
\item Supplied moderate-confidence OPEN-TRIAGE summary mentions Arnold's problems and continued-fraction literature without complete additional references or a verified resolution.
\end{itemize}
\end{document}
